\documentclass[10.5pt]{article}
\usepackage[letterpaper, margin=0.6in, top=0.5in, bottom=0.5in]{geometry}
\usepackage{amsmath,amssymb,amsthm}
\usepackage{fancyhdr}
\usepackage{xcolor}
\usepackage{tikz}
\usepackage{enumitem}
\usepackage{tcolorbox}
\usepackage{hyperref}

% Define colors matching instructor board annotations
\definecolor{chalkblue}{RGB}{0, 105, 185}
\definecolor{chalkred}{RGB}{205, 35, 35}
\definecolor{chalkpink}{RGB}{230, 90, 120}
\definecolor{chalkcyan}{RGB}{15, 160, 190}
\definecolor{chalkgreen}{RGB}{20, 130, 50}
\definecolor{boxbg}{RGB}{252, 252, 254}
\definecolor{boxframe}{RGB}{150, 160, 180}

\pagestyle{fancy}
\fancyhf{}
\renewcommand{\headrulewidth}{0pt}
\renewcommand{\footrulewidth}{0pt}

\setlength{\parindent}{0pt}
\setlength{\parskip}{2.5pt}

\begin{document}

% ==============================================================================
% PAGE 1
% ==============================================================================
\noindent \textbf{\large Math 105L} \hfill \textbf{\large Lab 2 Part 1}\\[1pt]
{\textbf{\Large Linear Models: Velocity \& Acceleration in Free-Fall}}\\[-5pt]
\rule{\textwidth}{1.2pt}

\vspace{0.03in}

\textbf{\large 1. Experimental Setup \& Data Cleaning}

An ultrasound motion detector records the height $h(t)$ (in meters) of a falling object every $0.03$ seconds:
\begin{itemize}[leftmargin=1.2em, topsep=1pt, itemsep=2pt]
  \item \textbf{Data Cleaning:} For $t > 0.42\text{ s}$, sensor reflections cause erroneous rebounds where the object falsely appears to rise. These points are removed to isolate true unhindered free-fall on $[0, 0.42]$.
  \item \textbf{Average Velocity Computation:} For each interval $[t_k, t_{k+1}]$, the average rate of change of height is:
  \[
    \mathbf{v_{\text{avg}} = \frac{h(t_{k+1}) - h(t_k)}{t_{k+1} - t_k} \quad \text{(units: meters/second)}}
  \]
  \item \textbf{Midpoint Assignment Rule:} To maximize numerical accuracy, the computed average velocity over $[t_k, t_{k+1}]$ is associated with the \textbf{center of the time interval}:
  \[
    \mathbf{t_{\text{mid}} = \frac{t_k + t_{k+1}}{2}}
  \]
\end{itemize}

\vspace{0.03in}

\textbf{\large 2. Classroom Demonstration on Blackboard: The First Time Step}

Over the interval from $t = 0$ to $t = 0.06$ seconds ($h(0) = 1.542\text{ m}$, $h(0.06) = 1.527\text{ m}$):
\begin{itemize}[leftmargin=1.2em, topsep=1pt, itemsep=3pt]
  \item \textbf{Midpoint Time:} $t_{\text{mid}} = \frac{0 + 0.06}{2} = \mathbf{0.03\text{ s}}$.
  \item \textbf{Average Velocity:}
  \[
    v_{\text{avg}} = \frac{1.527 - 1.542}{0.06 - 0} = \frac{-0.015}{0.06} = \mathbf{-0.25\text{ m/s}}
  \]
  \item Near $t = 0$, the height curve is almost horizontal, explaining why the initial velocity is small in magnitude.
\end{itemize}

\vspace{0.03in}

\begin{tcolorbox}[colback=boxbg, colframe=boxframe, boxrule=0.8pt, arc=3pt, left=8pt, right=8pt, top=4pt, bottom=4pt]
\textbf{\large The Velocity vs. Time Scatter Plot: A Linear Model}
When the calculated average velocities are plotted against the interval midpoints, the points align along a clear \textbf{straight line}:
\[
  \mathbf{v(t) = m t + b}
\]
\begin{itemize}[leftmargin=1.2em, topsep=1pt, itemsep=3pt]
  \item \textbf{Negative Velocities:} Indicate downward directional motion ($h(t)$ is decreasing).
  \item \textbf{Velocities Becoming More Negative:} The object is speeding up downward.
  \item \textbf{Slope of the Velocity Line ($m$):} Represents the rate of change of velocity:
  \[
    \mathbf{\text{Slope } m = \frac{\Delta v}{\Delta t} = \text{Acceleration } a \approx -9.8\text{ m/s}^2}
  \]
\end{itemize}
\end{tcolorbox}

\vspace{0.04in}

\textbf{\large 3. Core Synthesis}

\begin{tcolorbox}[colback=chalkblue!6, colframe=chalkblue, boxrule=1pt, arc=3pt, left=8pt, right=8pt, top=4pt, bottom=4pt]
\begin{itemize}[leftmargin=1.2em, topsep=1pt, itemsep=2pt]
  \item While the position graph $h(t)$ is \textbf{quadratic (curved downward)}, its derivative $v(t)$ is \textbf{linear}.
  \item The constant slope of this linear velocity model gives the physical constant of downward gravitational acceleration $g \approx 9.8\text{ m/s}^2$.
\end{itemize}
\end{tcolorbox}

\end{document}

